\chapter{Flight Dynamics and Control}
\label{chapter:papertwo}

\section{Coordinate Frames and System Definition}
\label{sec:coord_frames}

This chapter presents the flight dynamics and control framework for the
vehicle.  Its defining feature is that the airframe \emph{rotates continuously
about its vertical axis} while flying.  Two of the four rotors are mounted
horizontally, tangential to the frame, and drive this rotation directly at a
commanded thrust setpoint; the other two are mounted vertically and provide
lift and attitude authority.  All four turn in the same physical direction, so
the reaction torque of the lift rotors adds to the tangential drive rather than
cancelling it.  The vehicle therefore sustains a large, persistent yaw rate by
design---and, because the tangential rotors are commanded independently of the
lift channel, at a rate that is set rather than inherited from the thrust
balance.  This sustained rotation is what shapes the dynamics and motivates the
gyroscopic analysis that follows.

\paragraph{Inertial and body frames.}
The right-handed inertial frame $(\hat{x}^w, \hat{y}^w, \hat{z}^w)$ is
fixed to the ground with $\hat{z}^w$ pointing vertically upward.  The
position of the vehicle's centre of mass (COM) in this frame is
$\boldsymbol{p} = [x,\,y,\,z]^{\top}$.  The right-handed body frame
$(\hat{x}^b, \hat{y}^b, \hat{z}^b)$ is rigidly attached to the vehicle with
its origin at the COM; the $\hat{z}^b$ axis is aligned with the common
thrust axis of the four propellers.  The rotation matrix
$\boldsymbol{R}\in SO(3)$ maps body-frame vectors to the inertial frame,
so the world-frame direction of the thrust axis is $\boldsymbol{R}\hat{e}_3$,
where $\hat{e}_3 = [0,0,1]^{\top}$.  The body angular velocity expressed in
the body frame is $\boldsymbol{\omega} = [\omega_x, \omega_y, \omega_z]^{\top}$.
The yaw angle $\psi$ is defined as the rotation of $\hat{x}^b$ about
$\hat{z}^w$ relative to $\hat{x}^w$.

\paragraph{Non-revolving frame.}
Because the vehicle spins at high yaw rate, a non-revolving intermediate
frame $(\hat{x}^m, \hat{y}^m, \hat{z}^b)$ is introduced to describe the
attitude of the spinning disk independently of the instantaneous yaw
angle~\cite{bai2022splitflyerair}.  Under the small-inclination assumption
$|\alpha_x|, |\alpha_y| \ll 1$, the reduced attitude state
$\boldsymbol{\eta} = [\alpha_x, \alpha_y]^{\top}$ parametrises the deviation
of $\hat{z}^b$ from vertical, and the non-revolving axes are approximated as
\begin{equation}
    \hat{x}^m \approx [1,\;0,\;{-\alpha_y}]^{\top},
    \qquad
    \hat{y}^m \approx [0,\;1,\;\alpha_x]^{\top}.
    \label{eq:nm_frame}
\end{equation}
The pair $(\hat{x}^m, \hat{y}^m)$ spans the horizontal plane for any yaw
angle $\psi$; this frame is the natural setting for the gyroscopic reduced
attitude dynamics derived in Section~\ref{sec:reduced_dynamics}.

\paragraph{Rotor geometry and roles.}
The four rotors sit at equal arm length $l$ from the COM, at the body-frame
positions
\begin{equation}
    \boldsymbol{l}_1 = -\boldsymbol{l}_3
        = \frac{l}{\sqrt{2}}\!\left(\hat{x}^b + \hat{y}^b\right),
    \qquad
    \boldsymbol{l}_2 = -\boldsymbol{l}_4
        = \frac{l}{\sqrt{2}}\!\left(\hat{x}^b - \hat{y}^b\right),
    \label{eq:prop_locs}
\end{equation}
but they serve two distinct functions.  Rotors $1$ and $3$ are mounted with
their axes along $+\hat{z}^b$ and produce the \emph{lift thrusts}
$f_{1,p}, f_{3,p} \geq 0$ that support the vehicle and, through their
difference, tilt it.  Rotors $2$ and $4$ are mounted with their axes in the
body horizontal plane, oriented tangentially, and produce thrusts
$f_{2,p}, f_{4,p} \geq 0$ that act at moment arm $l$ about $\hat{z}^b$ and so
contribute only yaw torque.

The yaw torque therefore has two contributions in the same sense.  The
tangential rotors supply
\begin{equation}
    \tau_{z,\mathrm{tan}} = l\,(f_{2,p} + f_{4,p}),
    \label{eq:yaw_tangential}
\end{equation}
while each lift rotor contributes its aerodynamic reaction torque
$c\,f_{i,p}$ about $\hat{z}^b$, where $c > 0$ is the thrust-to-drag ratio
common to rotors sharing a spin direction.  Since the lift rotors turn in the
same direction as the tangential pair, these reaction torques add:
\begin{equation}
    \tau_z = \underbrace{l\,(f_{2,p} + f_{4,p})}_{\text{commanded}}
             + \underbrace{c\,(f_{1,p} + f_{3,p})}_{\text{reaction}}.
    \label{eq:yaw_torque}
\end{equation}
The first term dominates and is commanded directly: the tangential thrusts are
held at a fixed setpoint independent of the lift channel, so the spin rate is
selected rather than inherited.  The second term is a residual that varies
with the lift demand.  Because both terms drive the vehicle the same way, no
cancellation is available and no yaw equilibrium need be sought---the vehicle
spins because it is driven to, and $\omega_z$ is treated throughout this
chapter as a known, approximately constant operating parameter, measured in
flight and set by the choice of tangential thrust.

The total lift acting along $\hat{z}^b$ is carried by the vertical pair alone,
\begin{equation}
    F = f_{1,p} + f_{3,p},
    \label{eq:total_thrust_def}
\end{equation}
since the tangential rotors produce no component along the thrust axis.

\paragraph{Inertia and drag.}
By design, mass is distributed symmetrically about the $\hat{z}^b$ axis so
that the moments of inertia about roll and pitch are approximately equal,
$I_x \approx I_y \equiv I_d$.  The full inertia tensor is therefore
\begin{equation}
    \boldsymbol{I} = \mathrm{diag}(I_d,\;I_d,\;I_z),
    \label{eq:inertia}
\end{equation}
which is the axisymmetric condition required to treat the vehicle as a
gyroscope in Section~\ref{sec:reduced_dynamics}.  Rotor-induced drag is
characterised by the coefficient matrix
$\boldsymbol{B} = \mathrm{diag}(B_h, B_h, B_v)$, and aerodynamic body
damping by $\boldsymbol{D} = \mathrm{diag}(D_h, D_h, D_v)$, where subscripts
$h$ and $v$ denote horizontal and vertical components
respectively~\cite{bai2022splitflyerair}.  The vehicle mass is $m = 90\;\mathrm{g}$;
remaining hardware parameters are specified in Chapter~\ref{chapter:paperone}.

\paragraph{Propeller aerodynamic model.}
Each propeller $i$ spins about its own thrust axis at blade speed $\Omega_i$,
which is distinct from and substantially larger than the vehicle yaw rate
$\omega_z$.  The propulsive thrust follows the standard aerodynamic model
\begin{equation}
    f_{i,p} = k_f\,\Omega_i^2,
    \label{eq:thrust_model}
\end{equation}
where $k_f > 0$ is a propeller constant.  The hydrofoils are treated as
aerodynamically inert in air: they are thin flat plates carried at a fixed
inclination for water entry, not lifting surfaces, and the forces they generate
while sweeping through air are negligible beside the propeller thrust.  Their
aerodynamic role is confined to the brief submersion phase, where the working
fluid is water and the resulting forces are three orders of magnitude larger
(Chapter~\ref{dynamic}).

\section{Translational Dynamics}
\label{sec:translational_dynamics}

Treating the vehicle as a rigid body of mass $m$, Newton's second law gives
\begin{equation}
    m\ddot{\boldsymbol{p}} = \boldsymbol{R}\sum_{i=1}^{4}\boldsymbol{f}_i
        - mg\hat{e}_3,
    \label{eq:trans_eom}
\end{equation}
where the total force from propeller $i$ is
\begin{equation}
    \boldsymbol{f}_i = \hat{e}_3 f_{i,p} + \boldsymbol{f}_{i,d}.
    \label{eq:prop_force}
\end{equation}
The propulsive thrust $f_{i,p}$ is directed along the body $\hat{z}^b$ axis.
The rotor-induced drag $\boldsymbol{f}_{i,d}$, modelled as a linear function
of the air velocity at the propeller disk~\cite{bai2022splitflyerair}, is
\begin{equation}
    \boldsymbol{f}_{i,d} = -\boldsymbol{B}
        \!\left(\boldsymbol{R}^{-1}\dot{\boldsymbol{p}}
        + \boldsymbol{\omega}\times\boldsymbol{l}_i\right),
    \label{eq:rotor_drag}
\end{equation}
where the first term accounts for the translational velocity of the vehicle
and the second for the velocity of the propeller hub due to body rotation.

\paragraph{Cancellation of rotation-induced drag.}
The symmetry of the X-configuration---$\boldsymbol{l}_1 + \boldsymbol{l}_2
+ \boldsymbol{l}_3 + \boldsymbol{l}_4 = \mathbf{0}$, which follows
immediately from~\eqref{eq:prop_locs}---ensures that the rotation-induced
drag terms cancel when summed:
\begin{equation}
    \sum_{i=1}^{4}\left(\boldsymbol{\omega}\times\boldsymbol{l}_i\right)
        = \boldsymbol{\omega}\times\sum_{i=1}^{4}\boldsymbol{l}_i
        = \mathbf{0}.
    \label{eq:drag_cancel}
\end{equation}
This cancellation holds regardless of propeller spin direction and therefore
applies equally to our vehicle.  Substituting into~\eqref{eq:trans_eom}:
\begin{equation}
    m\ddot{\boldsymbol{p}}
        = \boldsymbol{R}\hat{e}_3 F
          - 4\boldsymbol{R}\boldsymbol{B}\boldsymbol{R}^{-1}\dot{\boldsymbol{p}}
          - mg\hat{e}_3.
    \label{eq:trans_full}
\end{equation}
The second term is a translational damping force present in all multirotor
vehicles.  For non-aggressive near-hover flight it is small relative to
gravity and is neglected in what follows~\cite{bai2022splitflyerair}, yielding
\begin{equation}
    m\ddot{\boldsymbol{p}} \approx \boldsymbol{R}\hat{e}_3 F - mg\hat{e}_3.
    \label{eq:trans_hover}
\end{equation}

\paragraph{Near-hover linearisation.}
For small inclinations $|\alpha_x|,|\alpha_y|\ll 1$, the body $\hat{z}^b$
axis expressed in the inertial frame is approximately
\begin{equation}
    \boldsymbol{R}\hat{e}_3 \approx
        \begin{bmatrix} J\boldsymbol{\eta} \\ 1 \end{bmatrix},
    \qquad
    J = \begin{bmatrix} 0 & 1 \\ -1 & 0 \end{bmatrix},
    \label{eq:zb_approx}
\end{equation}
where $J$ is the $-90^{\circ}$ rotation matrix in the horizontal plane.
Splitting~\eqref{eq:trans_hover} into horizontal and vertical components
and using the hover thrust condition $F\approx mg$ for the horizontal term
gives
\begin{align}
    m\ddot{\boldsymbol{p}}_2 &= mg\,J\boldsymbol{\eta},
    \label{eq:horiz_dyn}\\
    m\ddot{z} &= F - mg,
    \label{eq:vert_dyn}
\end{align}
where $\boldsymbol{p}_2 = [x,\,y]^{\top}$ is the horizontal position.
Equation~\eqref{eq:horiz_dyn} shows that horizontal acceleration is driven
solely by the vehicle tilt $\boldsymbol{\eta}$, and forms the basis of the
position controller in Section~\ref{sec:controller}.

\section{Attitude Dynamics}
\label{sec:attitude_dynamics}

The rotational motion is governed by Euler's equations in the body frame.
The angular momentum equation is
\begin{equation}
    \boldsymbol{I}\dot{\boldsymbol{\omega}}
        + \boldsymbol{\omega}\times\boldsymbol{I}\boldsymbol{\omega}
        = \sum_{i=1}^{4}\boldsymbol{\tau}_i
          + \boldsymbol{\tau}_s,
    \label{eq:euler}
\end{equation}
where $\boldsymbol{\tau}_i$ is the net torque produced by propeller $i$
at the centre of mass, and $\boldsymbol{\tau}_s = -\boldsymbol{D}\boldsymbol{\omega}$
is the aerodynamic body damping.

\paragraph{Propeller torques.}
The total torque from propeller $i$ decomposes into a thrust-induced
component and a drag-induced component:
\begin{equation}
    \boldsymbol{\tau}_i = \boldsymbol{\tau}_{i,p} + \boldsymbol{\tau}_{i,d}.
    \label{eq:prop_torque_decomp}
\end{equation}
For the two lift rotors ($i \in \{1,3\}$), the thrust-induced torque combines
the moment of the propulsive force about the centre of mass with the
aerodynamic reaction torque of the spinning
blade~\cite{bai2022splitflyerair}:
\begin{equation}
    \boldsymbol{\tau}_{i,p}
        = \boldsymbol{l}_i \times f_{i,p}\hat{e}_3
          + c\,f_{i,p}\hat{e}_3,
    \qquad i \in \{1,3\}.
    \label{eq:tau_thrust}
\end{equation}
The first term generates the differential roll/pitch torque used for attitude
control; the second is the reaction contribution to yaw.  For the two
tangential rotors ($i \in \{2,4\}$), whose thrust lies in the body horizontal
plane along $\hat{\boldsymbol{t}}_i$, the moment about the COM is directed
along $\hat{z}^b$,
\begin{equation}
    \boldsymbol{\tau}_{i,p}
        = \boldsymbol{l}_i \times f_{i,p}\hat{\boldsymbol{t}}_i
        = l\,f_{i,p}\,\hat{e}_3,
    \qquad i \in \{2,4\},
    \label{eq:tau_tangential}
\end{equation}
so these rotors contribute yaw torque alone and no roll or pitch.  Summing the
$\hat{z}^b$ components across all four rotors recovers~\eqref{eq:yaw_torque}.

The drag-induced torque arises from the rotor drag force
$\boldsymbol{f}_{i,d}$ (defined in~\eqref{eq:rotor_drag}) acting at the
moment arm $\boldsymbol{l}_i$:
\begin{equation}
    \boldsymbol{\tau}_{i,d}
        = \boldsymbol{l}_i \times \boldsymbol{f}_{i,d}
        = -\boldsymbol{l}_i
            \times \boldsymbol{B}
            \!\left(\boldsymbol{R}^{-1}\dot{\boldsymbol{p}}
                  + \boldsymbol{\omega}\times\boldsymbol{l}_i\right).
    \label{eq:tau_drag}
\end{equation}

\paragraph{Complete Euler equation.}
Collecting all torque contributions, with the propeller torques
$\boldsymbol{\tau}_{i,p}$ and $\boldsymbol{\tau}_{i,d}$ acting at each
rotor hub:
\begin{equation}
    \boldsymbol{I}\dot{\boldsymbol{\omega}}
        + \boldsymbol{\omega}\times\boldsymbol{I}\boldsymbol{\omega}
        = \sum_{i=1}^{4}
            \bigl(\boldsymbol{\tau}_{i,p}
                  + \boldsymbol{\tau}_{i,d}\bigr)
          + \boldsymbol{\tau}_s.
    \label{eq:euler_full}
\end{equation}

\paragraph{A note on the spinning hydrofoils.}
A rotating array of surfaces sweeping through air would, in principle, meet the
oncoming flow asymmetrically when the vehicle also translates, producing a
first-harmonic lift imbalance and hence a torque coupling translation into
tilt---the advancing/retreating blade effect familiar from rotorcraft.  This
coupling is not modelled here.  The hydrofoils are thin flat plates carried at
a fixed inclination for water entry rather than pitched lifting surfaces, and
the forces they generate in air are small enough to be absorbed into the
aerodynamic damping $\boldsymbol{\tau}_s$ already present
in~\eqref{eq:euler_full}.  The controller of Section~\ref{sec:controller}
accordingly makes no attempt to compensate such a term, and the flight results
of Chapter~\ref{chapter:paperthree} were obtained without it.  Quantifying this
coupling---which would require measuring the foils' aerodynamic characteristics
in air---is left to future work.

\paragraph{Yaw dynamics.}
Extracting the $\hat{z}^b$ projection of~\eqref{eq:euler_full} and using
$[\boldsymbol{\omega}\times\boldsymbol{I}\boldsymbol{\omega}]_z = 0$
(axisymmetric inertia) gives the yaw dynamics
\begin{equation}
    I_z\dot{\omega}_z
        = l\,(f_{2,p} + f_{4,p})
          + c\,(f_{1,p} + f_{3,p})
          - D_v\,\omega_z,
    \label{eq:yaw_ode}
\end{equation}
where the final term collects the aerodynamic resistance of the airframe to
rotation.  The driving terms are the two contributions
of~\eqref{eq:yaw_torque}; no attempt is made to model the resisting term in
detail, since the spin rate is not obtained by solving this balance.  The
tangential thrusts $f_{2,p}$ and $f_{4,p}$ are held at a fixed commanded
setpoint, so~\eqref{eq:yaw_ode} settles at whatever $\omega_z$ that setpoint
sustains, and a different setpoint yields a different spin rate.  Raising the
commanded tangential thrust raises $\omega_z$; the spin rate is thus an
operating parameter chosen by the pilot or the experiment, and is measured in
flight rather than predicted.

\paragraph{The spin rate as a design variable.}
This is a substantive departure from vehicles that spin because their yaw
torque cannot be cancelled.  In those platforms the spin rate is whatever the
airframe settles at, and the design problem is to arrange for that value to be
tolerable~\cite{mueller2016relaxed,bai2020splitflyer}.  Here the dominant yaw
drive is an independent input, so $\omega_z$ can be selected---and the relevant
question becomes what value to select.

The choice is governed from outside the flight dynamics.  The hydrofoils
sweep at tangential speed $\omega_z r_f$, and the hydrodynamic force they
generate on water contact grows with the square of that speed
(Chapter~\ref{dynamic}); the spin rate is therefore set by what the hop
requires, subject to retaining enough spin through the contact for the vehicle
to remain controllable on exit.  In flight testing the vehicle was operated at
spin rates approaching $10{,}000^{\circ}\,\mathrm{s^{-1}}$ with no
controllability limit encountered, so this upper end of the range was not
found to be restrictive in practice.

What the analysis that follows requires of $\omega_z$ is only that it be large
and approximately constant over the timescale of the attitude dynamics.  Both
hold: the tangential thrust setpoint is fixed during flight, and the resulting
spin is slow to change because $I_z$ is large relative to the available yaw
torque.  The value of $\omega_z$ used in the reduced model of
Section~\ref{sec:reduced_dynamics} is the measured one.

\section{Reduced Dynamics under High Yaw Rate}
\label{sec:reduced_dynamics}

The full Euler equation~\eqref{eq:euler_full} governs all three rotational
degrees of freedom simultaneously.  The equilibrium analysis of
Section~\ref{sec:attitude_dynamics} established that the vehicle operates at
an approximately constant yaw rate $\omega_z$ set by the commanded tangential
thrust, with the tilt rates $|\omega_x|, |\omega_y|$ remaining small compared
to $\omega_z$.  This separation of timescales motivates a reduced
two-dimensional description in the tilt state $\boldsymbol{\eta}$ that
eliminates the fast yaw and retains only the controllable attitude evolution.

\paragraph{Inertial angular momentum.}
The angular momentum in the inertial frame is
\begin{equation}
    \mathbf{h} = \boldsymbol{R}\boldsymbol{I}\boldsymbol{\omega}
        = I_d\boldsymbol{R}[\omega_x,\;\omega_y,\;0]^{\top}
          + I_z\omega_z\boldsymbol{R}\hat{e}_3.
    \label{eq:inertial_h}
\end{equation}
The second term is the large gyroscopic spin momentum $I_z\omega_z$
directed along the tilting thrust axis $\boldsymbol{R}\hat{e}_3$.
By Newton's second law in the inertial frame,
\begin{equation}
    \dot{\mathbf{h}}
        = \boldsymbol{R}\!\left(
            \sum_{i=1}^{4}\bigl(\boldsymbol{\tau}_{i,p}
                               + \boldsymbol{\tau}_{i,d}\bigr)
            + \boldsymbol{\tau}_{AB}
            + \boldsymbol{\tau}_s
          \right).
    \label{eq:h_dot}
\end{equation}

\paragraph{Horizontal projection and gyroscopic coupling.}
Projecting~\eqref{eq:h_dot} onto the inertial horizontal plane via
$[e_1,e_2]^{\top}$ and using $[\omega_x, \omega_y]^{\top} \approx
\dot{\boldsymbol{\eta}}$ for small tilt, the two parts
of~\eqref{eq:inertial_h} contribute separately.  The tilt momentum term
gives $I_d\ddot{\boldsymbol{\eta}}$ on differentiation.  The spin momentum
term $I_z\omega_z\boldsymbol{R}\hat{e}_3$ precesses as the body tilts:
\begin{equation}
    [e_1,e_2]^{\top}
    \frac{\mathrm{d}}{\mathrm{d}t}
    \bigl(I_z\omega_z\boldsymbol{R}\hat{e}_3\bigr)
    = I_z\omega_z[e_1,e_2]^{\top}
      \boldsymbol{R}\bigl(\boldsymbol{\omega}\times\hat{e}_3\bigr)
    \approx I_z\omega_z\boldsymbol{J}\dot{\boldsymbol{\eta}},
    \label{eq:gyro_proj}
\end{equation}
where the last step uses $\boldsymbol{\omega}\times\hat{e}_3 \approx
\boldsymbol{J}\boldsymbol{\omega}_h = [\omega_y,-\omega_x]^{\top}$ and the
fact that $\boldsymbol{J}$ commutes with planar rotations (since $\boldsymbol{J}$
is itself a $-90^{\circ}$ rotation in the horizontal
plane)~\cite{bai2022splitflyerair}.  The result is the \emph{gyroscopic
coupling} term $I_z\omega_z\boldsymbol{J}\dot{\boldsymbol{\eta}}$: the
large spin angular momentum precesses in the inertial plane as the vehicle
tilts, coupling $\dot{\alpha}_x$ into the $\alpha_y$ equation and vice
versa.

\paragraph{Projected external torques.}
Two torque groups appear on the right side of~\eqref{eq:h_dot}.

\noindent\textit{Control torques.}
The inertial projection of the thrust-induced roll/pitch torques defines the
\emph{reduced control input}:
\begin{equation}
    \tilde{\boldsymbol{\tau}}_w
        = [e_1,e_2]^{\top}\boldsymbol{R}
          \sum_{i=1}^{4}\bigl(\boldsymbol{l}_i\times f_{i,p}\hat{e}_3\bigr).
    \label{eq:tau_w_reduced}
\end{equation}
The yaw reaction torques $cf_{i,p}\hat{e}_3$ project to $cF\boldsymbol{J}\boldsymbol{\eta}$,
which is second-order in tilt amplitude and is neglected.

\noindent\textit{Drag torques.}
The horizontal arm-sweep component of the rotor drag produces only yaw
torque (its roll/pitch contribution cancels by
symmetry~\eqref{eq:drag_cancel}).  The vertical rotor drag at each hub,
arising from the tilt-induced hub velocity $(\boldsymbol{\omega}\times
\boldsymbol{l}_i)_z$, contributes a roll/pitch damping; summing over four
arms of equal length gives $-2B_v l^2\dot{\boldsymbol{\eta}}$.  Adding the
body damping $-D_h\dot{\boldsymbol{\eta}}$ from $\boldsymbol{\tau}_s$ defines
the \emph{effective tilt damping coefficient}
\begin{equation}
    \bar{D} \triangleq 2B_v l^2 + D_h,
    \label{eq:Dbar}
\end{equation}
so the projected drag torque is $\tilde{\boldsymbol{\tau}}_d =
-\bar{D}\dot{\boldsymbol{\eta}}$.

\paragraph{Reduced attitude equation.}
Combining the projected angular momentum with the external torques yields
the \emph{gyroscopic reduced attitude equation}:
\begin{equation}
    \boxed{I_d\ddot{\boldsymbol{\eta}}
        + I_z\omega_z\boldsymbol{J}\dot{\boldsymbol{\eta}}
        = \tilde{\boldsymbol{\tau}}_w
          - \bar{D}\dot{\boldsymbol{\eta}}.}
    \label{eq:reduced_attitude}
\end{equation}
The gyroscopic term $I_z\omega_z\boldsymbol{J}\dot{\boldsymbol{\eta}}$
couples the two tilt channels through precession of the spin angular momentum,
and $-\bar{D}\dot{\boldsymbol{\eta}}$ is the projected aerodynamic damping.

Together with the horizontal position equation~\eqref{eq:horiz_dyn}:
\begin{equation}
    m\ddot{\boldsymbol{p}}_2 = mg\boldsymbol{J}\boldsymbol{\eta},
    \label{eq:horiz_coupled}
\end{equation}
equations~\eqref{eq:reduced_attitude}--\eqref{eq:horiz_coupled} form a
coupled fourth-order system in $(\boldsymbol{p}_2,\boldsymbol{\eta})$ driven
by $\tilde{\boldsymbol{\tau}}_w$.  This is the plant model for the controller
developed in Section~\ref{sec:controller}.

\paragraph{Actuation: degrees of freedom and thrust generation.}
The reduced input $\tilde{\boldsymbol{\tau}}_w$ is a two-component vector, and
the actuation available to realise it differs from that of a conventional
quadrotor because the four rotors are not interchangeable.  The tangential
pair is held at its commanded spin setpoint and is not modulated for attitude
control, leaving the two lift rotors to supply both the total lift $F$ and the
tilting torque.  Their thrusts map to the lift and body-frame torque as
\begin{equation}
    \begin{bmatrix}F\\\tau^b\end{bmatrix}
    = \boldsymbol{M}
      \begin{bmatrix}f_{1,p}\\f_{3,p}\end{bmatrix},
    \qquad
    \boldsymbol{M} = \begin{bmatrix}
        1 & 1 \\[2pt]
        l & -l
    \end{bmatrix},
    \label{eq:input_map}
\end{equation}
obtained from~\eqref{eq:tau_thrust} with the rotor
positions~\eqref{eq:prop_locs}.  The sum of the two thrusts sets the lift and
their difference sets the torque, so $\boldsymbol{M}$ is square and invertible:
a given $(F,\tau^b)$ determines the two lift thrusts uniquely, subject to the
positivity constraints $f_{i,p}\geq 0$.

The two lift rotors provide differential thrust about a single axis fixed in
the body frame, whereas the required $\tilde{\boldsymbol{\tau}}_w$ is a
two-dimensional vector in the inertial plane.  The vehicle's rotation
reconciles the two: as the airframe spins, that body-fixed axis sweeps through
every inertial direction once per revolution, so a differential thrust whose
magnitude and sign vary with the yaw angle produces a torque of any desired
inertial direction~\cite{bai2020splitflyer,bai2022splitflyerair}.  The
required modulation is exactly the projection of the inertial command onto the
instantaneous body axes, given by~\eqref{eq:torque_projection} below and
evaluated at the measured yaw angle on every control cycle.

\section{Controller Design}
\label{sec:controller}

The control objective is to track a desired horizontal trajectory
$\boldsymbol{p}_{2,d}(t)$ and altitude set-point $z_d$ while the vehicle spins
at the commanded rate of Section~\ref{sec:attitude_dynamics}.  The spin itself
is not regulated in closed loop: the tangential thrust is held at its setpoint
and the resulting $\omega_z$ is measured and used by the controller, but no
feedback acts on it.

\paragraph{Non-cascaded horizontal controller.}
Define the horizontal position error $\tilde{\boldsymbol{e}} =
\boldsymbol{p}_2 - \boldsymbol{p}_{2,d}$.  From~\eqref{eq:horiz_coupled},
the attitude and its rate are related to position derivatives by
\begin{equation}
    \boldsymbol{\eta} = -\frac{1}{g}\boldsymbol{J}\ddot{\boldsymbol{p}}_2,
    \qquad
    \dot{\boldsymbol{\eta}} = -\frac{1}{g}\boldsymbol{J}\boldsymbol{p}_2^{(3)},
    \label{eq:eta_from_p2}
\end{equation}
so the system~\eqref{eq:reduced_attitude}--\eqref{eq:horiz_coupled} is fourth
order in $\boldsymbol{p}_2$.  Solving~\eqref{eq:reduced_attitude} for
$\tilde{\boldsymbol{\tau}}_w$ yields the plant input--output relation:
\begin{equation}
    \tilde{\boldsymbol{\tau}}_w
        = -\frac{I_d}{g}\boldsymbol{J}\boldsymbol{p}_2^{(4)}
          + \bar{D}\dot{\boldsymbol{\eta}}
          + I_z\omega_z\boldsymbol{J}\dot{\boldsymbol{\eta}}.
    \label{eq:ctrl_inversion}
\end{equation}
The last two terms are the known plant dynamics---aerodynamic damping and
gyroscopic coupling---to be cancelled by the control law.

A fourth-order polynomial constraint is imposed on the error:
\begin{equation}
    \tilde{\boldsymbol{e}}^{(4)}
        + \lambda_3\tilde{\boldsymbol{e}}^{(3)}
        + \lambda_2\ddot{\tilde{\boldsymbol{e}}}
        + \lambda_1\dot{\tilde{\boldsymbol{e}}}
        + \lambda_0\tilde{\boldsymbol{e}} = \boldsymbol{0},
    \label{eq:error_poly}
\end{equation}
which is asymptotically stable when $s^4 + \lambda_3 s^3 + \lambda_2 s^2
+ \lambda_1 s + \lambda_0$ is Hurwitz~\cite{bai2022splitflyerair}.  The
constraint prescribes the reference fourth derivative:
\begin{equation}
    \boldsymbol{p}_2^{(4)}\big|_r
        = \boldsymbol{p}_{2,d}^{(4)}
          - \lambda_3\tilde{\boldsymbol{e}}^{(3)}
          - \lambda_2\ddot{\tilde{\boldsymbol{e}}}
          - \lambda_1\dot{\tilde{\boldsymbol{e}}}
          - \lambda_0\tilde{\boldsymbol{e}}.
    \label{eq:p2_ref}
\end{equation}
Substituting~\eqref{eq:p2_ref} into~\eqref{eq:ctrl_inversion} gives the
horizontal control torque:
\begin{equation}
    \tilde{\boldsymbol{\tau}}_{w,r}
        = -\frac{I_d}{g}\boldsymbol{J}\boldsymbol{p}_2^{(4)}\big|_r
          + \bar{D}\dot{\boldsymbol{\eta}}
          + I_z\omega_z\boldsymbol{J}\dot{\boldsymbol{\eta}}.
    \label{eq:ctrl_law}
\end{equation}
The first term is the error-driven command from~\eqref{eq:p2_ref}; the
remaining terms cancel the known plant dynamics.  The jerk
$\tilde{\boldsymbol{e}}^{(3)}$ in~\eqref{eq:p2_ref} is evaluated as
$g\boldsymbol{J}\dot{\boldsymbol{\eta}} - \boldsymbol{p}_{2,d}^{(3)}$
from~\eqref{eq:eta_from_p2}, so all feedback signals reduce to
$(\boldsymbol{p}_2,\,\dot{\boldsymbol{p}}_2,\,\boldsymbol{\eta},\,
\dot{\boldsymbol{\eta}})$---quantities available from position sensing and
an IMU.

\paragraph{Altitude controller.}
The vertical dynamics~\eqref{eq:vert_dyn} are controlled independently by
the total thrust.  A proportional-derivative law with gravity compensation
gives
\begin{equation}
    F_r = m\bigl[
        \ddot{z}_d
        + k_{v,z}(\dot{z}_d - \dot{z})
        + k_{p,z}(z_d - z)
    \bigr] + mg,
    \label{eq:alt_ctrl}
\end{equation}
where $k_{p,z},\,k_{v,z} > 0$ are the altitude proportional and derivative
gains.

\paragraph{Input allocation.}
The reference torque $\tilde{\boldsymbol{\tau}}_{w,r} = [u_x, u_y]^{\top}$ is
expressed in the inertial horizontal frame, while the lift rotors act about
axes that rotate with the vehicle.  The command is therefore resolved into the
rotating frame using the measured yaw angle $\psi$,
\begin{equation}
    \begin{bmatrix}\tau^b_{x,r}\\[2pt]\tau^b_{y,r}\end{bmatrix}
        = \boldsymbol{R}_{z,2}^{\top}(\psi)
          \begin{bmatrix}u_x\\[2pt]u_y\end{bmatrix}
        = \begin{bmatrix}
            -u_x\sin\psi + u_y\cos\psi \\[2pt]
             u_x\cos\psi + u_y\sin\psi
          \end{bmatrix},
    \label{eq:torque_projection}
\end{equation}
where $\boldsymbol{R}_{z,2}(\psi)$ is the $2\times2$ yaw-rotation submatrix of
$\boldsymbol{R}$.  Because $\psi$ advances continuously at the spin rate, this
projection modulates the body-frame command once per revolution, which is what
allows a body-fixed actuation axis to realise an arbitrary inertial torque.
The projection is evaluated on board at the attitude-loop rate from the yaw
angle supplied by the position-tracking system, so its phase remains
synchronised with the vehicle's rotation without requiring the outer loop to
run at the spin frequency.

The resolved commands then set the differential thrust of the lift pair,
\begin{equation}
    f_{1,r} = \frac{F_r}{2} + \frac{\tau^b_r}{2l},
    \qquad
    f_{3,r} = \frac{F_r}{2} - \frac{\tau^b_r}{2l},
    \qquad
    \tau^b_r = \tau^b_{y,r} - \tau^b_{x,r},
    \label{eq:alloc}
\end{equation}
with $F_r$ from~\eqref{eq:alt_ctrl}, the two rotors sharing the collective
lift and splitting the tilting torque between them.  The differential term is
limited so that both thrusts remain non-negative and within the motor
saturation bound, which caps the attainable tilting torque at low collective
thrust.

The tangential rotors take no part in this allocation.  Holding them at a
fixed setpoint independent of $F_r$ and $\tilde{\boldsymbol{\tau}}_{w,r}$
means the spin is sustained by construction whatever the lift channel is
doing---including when the lift command is removed entirely.  This is what
makes the descent of Chapter~\ref{chapter:paperthree} possible: the two lift
rotors are commanded to zero thrust and the vehicle falls ballistically onto
the water, while the tangential pair continues to drive the rotation that the
hydrofoils need at contact.  Under the coupled yaw--thrust arrangement of a
conventional spinning multirotor, cutting the lift would also remove the yaw
drive and the spin would decay exactly when it is needed most.